% !TeX program = xelatex
% Faithful UTF-8 transcription of all six handwritten source images.
% Medical statements are preserved as written, without editorial correction.
% pic01.jpg = 20261019_082357179.jpg
% pic02.jpg = 20261019_082404053.jpg
% pic03.jpg = 20261019_082411117.jpg
% pic04.jpg = 20261019_082417606.jpg
% pic05.jpg = 20261019_082426611.jpg
% pic06.jpg = 20261019_082433180.jpg
% Compile with XeLaTeX; also compatible with input into a XePersian document.
\ifdefined\settextfont
  \def\EGAStandalone{0}
\else
  \def\EGAStandalone{1}
  \documentclass[11pt,a4paper]{article}
  \usepackage[margin=1.8cm]{geometry}
  \usepackage{amsmath,amssymb,enumitem}
  \usepackage{xepersian}
  \IfFontExistsTF{Amiri}{\settextfont{Amiri}}{\settextfont{FreeSerif}}
  \setlatintextfont{FreeSerif}
  \setlength{\parindent}{0pt}
\fi
\providecommand{\en}[1]{\lr{#1}}
\newcommand{\EGAarrow}{\en{$\longleftarrow$}}
\newcommand{\EGAheader}{%
  {\footnotesize\en{Note book\qquad Subject: 16\qquad
  Veor: \_\_\_\quad Month: \_\_\_\quad Date: \_\_\_}\par}}
\newcommand{\EGAfooter}[1]{%
  \par\medskip{\footnotesize\en{BAHAR DANESH ASL}\hfill #1\par}}
\ifnum\EGAStandalone=1 \begin{document}\fi

% ---------------- pic01.jpg ----------------
\EGAheader
\section*{\en{Extraluminal gas in the abdomen}}
\begin{enumerate}[label=\arabic*.,leftmargin=*,itemsep=0.6em]
\item شایع‌ترین مکان‌ها اکسترالومینال \en{air}
  \begin{itemize}[nosep]
  \item اینتراپریتونئال (نوموپریتونئوم) \textbf{\en{(free air)}}
  \item رتروپریتونئال
  \item در دیوارهٔ روده \en{(pneumatosis intestinalis)}
  \item مجاری صفراوی \en{(pneumobilia)}
  \end{itemize}
\item \en{Radiographic Signs of free intraperitoneal air}
  \begin{latin}
  \begin{enumerate}[label=\alph*),leftmargin=*,nosep]
  \item Air beneath the diaphragm
  \item Visualization of both sides of the bowel wall
  \item Visualization of the falciform ligament
  \end{enumerate}
  \end{latin}
\item هوا زیر دیافراگم
  \begin{itemize}[itemsep=0.4em]
  \item \en{Crescent lucency}
  \item می‌تواند بعد از جراحی‌ها شکمی تا \textbf{۷ روز} دیده شود ولی در گرافی‌های سریال باید کاهش یابد.
  \item \textbf{\en{CT} حساس‌ترین روش بررسی هوا آزاد است.}
  \item \textbf{گرافی به‌عنوان یک ابزار غربالگری است.}
  \item تشخیص آن در سمت راست راحت‌تر است.
  \item اگر بیمار نتواند گرافی \en{Upright} بگیرد از گرافی \en{left lateral decubitus} استفاده می‌کنیم.
  \end{itemize}
\end{enumerate}
\EGAfooter{۱}

% ---------------- pic02.jpg ----------------
\clearpage
\EGAheader
\begin{enumerate}[label=\arabic*.,leftmargin=*,itemsep=0.7em]
\setcounter{enumi}{3}
\item \en{Chilaiditi Sign} (کیلاایتی)
  \begin{itemize}[itemsep=0.4em]
  \item در گرافی ایستاده، در برخی افراد کولون می‌تواند بین کبد و همی‌دیافراگم راست دیده شود که ممکن است با هوای آزاد اشتباه گرفته شود.
  \item \textbf{کولون \en{haustral fold} دارد.}
  \item اگر شک داشتیم \textbf{\en{LLD}} می‌گیریم و در مرحلهٔ بعد \en{CT}.
  \end{itemize}
\item در گرافی نرمال شکم ما فقط هوای داخل لومن را می‌بینیم و خود دیوارهٔ روده دیده نمی‌شود.
  \begin{itemize}[itemsep=0.4em]
  \item ورود هوا به حفرهٔ پریتوئن باعث می‌شود که بتوانیم دیوارهٔ روده را نیز ببینیم.
    \EGAarrow\ \textbf{\en{Rigler Sign}}
  \item به تمامی نماها دیده می‌شود.
  \end{itemize}
\item وقتی لوپ‌ها دیلاتهٔ روده کوچک روی هم می‌افتند ممکن است ظاهری شبیه به \en{Rigler Sign} ایجاد کنند.
  \begin{itemize}[itemsep=0.4em]
  \item \en{Pseudo-Rigler Sign}
  \item با گرافی یا \en{CT} اطمینان حاصل می‌کنیم.
  \end{itemize}
\item لیگامان فالسیفورم در لبهٔ آزاد کبد در قدام، در سمت راست \en{Upper lumbar Spine} قرار دارد.
  \begin{itemize}[nosep]
  \item به‌صورت طبیعی قابل مشاهده نیست.
  \end{itemize}
\end{enumerate}
\EGAfooter{۲}

% ---------------- pic03.jpg ----------------
\clearpage
\EGAheader
\begin{itemize}[itemsep=0.4em]
\item اگر هوای آزاد زیادی در شکم باشد \EGAarrow\ در گرافی سوپاین
  \begin{itemize}[nosep]
  \item هوای آزاد اطراف لیگامان فالسیفورم را می‌گیرد.
  \item لیگامان نمایان می‌شود.
  \item \en{Falciform Ligament Sign}
  \end{itemize}
\end{itemize}
\begin{enumerate}[label=\arabic*.,leftmargin=*,itemsep=0.6em]
\setcounter{enumi}{7}
\item ظاهر خطی لیگامان فالسیفورم در کنار ظاهر اووال تجمع هوا آزاد
  \EGAarrow\ \textbf{\en{Football Sign}}
\end{enumerate}
\subsection*{علل هوا آزاد}
\begin{enumerate}[label=\arabic*.,leftmargin=*,itemsep=0.6em]
\setcounter{enumi}{8}
\item شایع‌ترین علت \EGAarrow\ پارگی بخشی حاوی هوا در لولهٔ گوارش
  \begin{itemize}[nosep]
  \item معده
  \item \en{SB}
  \item \en{LB}
  \item از بین همه پپتیک اولسر پرفوره از همه شایع‌تر است.
  \end{itemize}
\item هوا آزاد تا ۷ روز بعد از جراحی شکمی طبیعی است.
  \begin{itemize}[nosep]
  \item اگر بیش از یک هفته وجود داشت
  \item به طول زمان مقدارش زیاد شد
  \end{itemize}
  \EGAarrow\ کامپلیکیشن جراحی یا بیماری اولیه
\item \textbf{دایورتیکولیت و آپاندیسیت پرفوره عموماً \en{walled-off abscess collections} ایجاد می‌کنند.}
  \begin{itemize}[nosep]
  \item هوای آزاد در نتیجهٔ آن‌ها ناشایع است.
  \end{itemize}
\item پرفوریشن ناشی از کارسینوم کولون به‌طور ناشایع می‌تواند هوا آزاد ایجاد نماید.
\end{enumerate}
\EGAfooter{۳}

% ---------------- pic04.jpg ----------------
\clearpage
\EGAheader
\section*{هوای رتروپریتونئال}
\begin{enumerate}[label=\arabic*.,leftmargin=*,itemsep=0.6em]
\setcounter{enumi}{12}
\item \en{Recognized by:}
  \begin{latin}
  \begin{enumerate}[label=\alph*),leftmargin=*,itemsep=0.4em]
  \item Streaky, linear appearance outlining extraperitoneal structures
  \item Mottled, blotchy appearance
    \rl{(خصوصاً فضای آنتریور پارارنال)}
  \item Relatively fixed positions
  \end{enumerate}
  \end{latin}
\item هوا رتروپریتونئال
  \begin{itemize}[itemsep=0.4em]
  \item از طریق دیافراگماتیک هیاتوس وارد مدیاستن شود
    \EGAarrow\ نومومدیاستن
  \item وارد پریتوئن شود \EGAarrow\ نوموپریتوئن
  \end{itemize}
\item علل هوای رتروپریتونئال
  \begin{itemize}[nosep]
  \item عموماً ناشی از پرفوریشن ثانویه به:
    \begin{itemize}[nosep]
    \item بیماری التهابی (آپاندیسیت پرفوره)
    \item بیماری اولسراتیو (کرون)
    \end{itemize}
  \end{itemize}
\end{enumerate}
\section*{نوماتوس اینتستینالیس}
\begin{enumerate}[label=\arabic*.,leftmargin=*,itemsep=0.5em]
\item عموماً در گرافی سادهٔ شکم قابل تشخیص است.
  \begin{latin}
  It is seen in profile producing a linear radiolucency (black line)
  whose contour exactly parallels the bowel lumen.
  \end{latin}
\end{enumerate}
\EGAfooter{۴}

% ---------------- pic05.jpg ----------------
\clearpage
\EGAheader
\begin{enumerate}[label=\arabic*.,leftmargin=*,itemsep=0.6em]
\setcounter{enumi}{1}
\item افتراق نوماتوس اینتستینالیس از \en{fecal material}
  \begin{itemize}[itemsep=0.4em]
  \item وجود \en{mottled gas} در نقطه‌ای که بعید است کولون باشد
  \item تغییر نکردن مکان آن در گرافی سریال یا تغییر پوزیشن
  \end{itemize}
\item علل نوماتوس اینتستینالیس
  \begin{itemize}[itemsep=0.4em]
  \item \en{Pneumatosis cystoides intestinalis} \EGAarrow\ نادر است.
    \begin{itemize}[nosep]
    \item فرم اولیهٔ نوماتوس اینتستینالیس است.
    \item \textbf{عموماً کولون چپ را درگیر می‌کند.}
    \item \en{Cystlike collections} \EGAarrow\ در سروز یا ساب‌موکوز
    \end{itemize}
  \item فرم شایع‌تر در پی \en{COPD} \EGAarrow\ فرم ثانویه
    \begin{itemize}[nosep]
    \item هوا از مدیاستن وارد شکم می‌شود.
    \end{itemize}
  \item سایر انواع فرم ثانویه
    \begin{itemize}[itemsep=0.3em]
    \item \textbf{\en{Necrotizing enterocolitis} در اطفال}
    \item \en{Ischemic bowel disease} (بزرگسالان)
    \item \en{Obstructing lesions of the bowel}
      \begin{itemize}[nosep]
      \item \en{Obstructing carcinomas in adults}
      \item \en{Hirschsprung or Pyloric stenosis In children}
      \end{itemize}
    \end{itemize}
  \end{itemize}
\item عوارض نوماتوس اینتستینالیس
  \begin{itemize}[itemsep=0.4em]
  \item \en{Rupture into the peritoneal cavity} \EGAarrow\ نوموپریتوئن
  \item \en{Dissection of air into the portal venous system}
  \end{itemize}
\end{enumerate}
\EGAfooter{۵}

% ---------------- pic06.jpg ----------------
\clearpage
\EGAheader
\section*{هوا در سیستم بیلیاری (نوموبیلیا)}
\begin{enumerate}[label=\arabic*.,leftmargin=*,itemsep=0.6em]
\setcounter{enumi}{4}
\item \en{Recognizing}
  \begin{latin}
  One or two tube-like, branching lucencies in the RUQ overlying the
  \textbf{central portion of the liver} that conform to the location
  and appearance of the major bile ducts.
  \end{latin}
\item \en{Portal Venous Gas}
  \begin{latin}
  \begin{itemize}[leftmargin=*,itemsep=0.3em]
  \item Most often associated with \textbf{necrotizing enterocolitis} in infants.
  \item Can be seen in adults usually with bowel necrosis.
  \item Small black branching structures over the \textbf{periphery} of the liver.
  \end{itemize}
  \end{latin}
\item \en{Signs of pneumobilia}
  \begin{latin}
  \begin{enumerate}[label=\alph*),leftmargin=*,itemsep=0.3em]
  \item Tube-like, branching lucencies in the RUQ
  \item Tubular structures $\longrightarrow$ Central in location
  \item Air may be seen in the lumen of gallbladder
  \end{enumerate}
  \end{latin}
\item علل نوموبیلیا
  \begin{itemize}[itemsep=0.4em]
  \item در صورت عدم کفایت اسفنکتر اودی یک یافتهٔ نرمال است.
  \item جراحی در مجاری صفراوی
  \item علل پاتولوژیک
    \begin{itemize}[nosep]
    \item \en{Gallstone ileus}
    \item \en{Gas forming pyogenic colangitis} \EGAarrow\ \en{E. Coli}
    \end{itemize}
  \end{itemize}
\end{enumerate}
\EGAfooter{۶}

\ifnum\EGAStandalone=1 \end{document}\fi
